% !PRINTED 192
\weberpage{192}
%<<CONT
Factoren soll $q$ betragen. Es ist dann, da jeder Factor zwei
Indices enth\"alt,
\weq{19}{\nu + n\mu = 2q.}
Dann ergiebt sich aus (17) und (18), wenn $P'$ aus $P$ durch
Vertauschung von $x, \alpha_i$ mit $x', \alpha_i'$ hervorgeht
\weq{20}{{a_0'}^{\mu} P' = a_0^{\mu} r^{n\mu-q} (\gamma x' + \delta)^{\nu} P.}
Wir bilden nun
\weq{21}{C(x, y, a) = a_0^{\mu} y^{\nu} \Sigma P\left(\frac{x}{y}\right),}
die Summe genommen \"uber alle Werthe, die man aus $P$ erh\"alt,
wenn man die Indices $1, 2, 3 \ldots n$ auf alle m\"ogliche Arten
mit einander vertauscht. Diese Summe ist dann eine sym\-
metrische Function der Wurzeln von (1) und l\"asst sich also als
ganze rationale Function der Verh\"altnisse $a_1 : a_0, a_2 : a_0 \ldots a_n : a_0$
ausdr\"ucken, so dass kein Glied die $\mu^{\text{te}}$ Ordnung \"ubersteigt. Die
Function $C(x, y, a)$ ist also eine ganze homogene Function von
$a_0, a_1 \ldots a_n$ von der $\mu^{\text{ten}}$ Ordnung und eine ganze homogene
Function von $x, y$ von der $\nu^{\text{ten}}$ Ordnung. Da sie nach (19), (20)
der Bedingung
% [sic] Im Druck fehlt in C(x', y' a') das Komma zwischen y' und a'.
\weq{22}{C(x', y' a') = r^{\lambda} C(x, y, a),}
\weq{23}{\lambda = n\mu - q, \quad 2\lambda = n\mu - \nu}
gen\"ugt, so ist sie eine \gesperrt{Covariante} und in dem besonderen
Falle, wo $\nu = 0$ ist, eine \gesperrt{Invariante}.

Das Gewicht $\lambda$ ist, wie der zweite Ausdruck zeigt, immer
positiv oder Null, da $n\mu$ mindestens gleich $\nu$. Der \"ausserste
Werth $n\mu = \nu$ oder $\lambda = 0$ kommt nur dann vor, wenn in $P$
jeder der Indices $1, 2, 3 \ldots n$ nur mit 0 verbunden vorkommt,
also $P$ eine Potenz von $f(x)$ selbst ist.

\begin{center}
\S. 62.

\gesperrt{Bin\"are cubische Formen.}
\end{center}

Eine bin\"are quadratische Form giebt keine anderen invarianten
 Bildungen als die Discriminante. Mit diesen befassen
wir uns daher nicht und gehen gleich zur Betrachtung der
cubischen Formen
\weq{1}{f(x, y) = a_0 x^3 + a_1 x^2 y + a_2 x y^2 + a_3 y^3}
%CONT>>
